

\def\topic{Evaluation}
\ifdefined\MyHeaderIncluded
\else
\def\qwertyflag{true}

\ifdefined\MyHeaderIncluded

\else
  \newcommand{\foobar}{FooBar}
  \def\MyHeaderIncluded{true}
% if is till the end of the file

\documentclass[10pt,a4paper]{article}
\usepackage[utf8]{inputenc}
%\usepackage[german]{babel}
%\usepackage[cyrillic]{babel}
\usepackage[T1]{fontenc}
%\usepackage[T2A]{fontenc}
%\usepackage[english, russian]{babel}
\usepackage{amsmath}
\usepackage{amsfonts}
\usepackage{amssymb}
\usepackage{amsthm}
\usepackage{mathtools} %coloneqq

\usepackage{bussproofs} %

%\renewcommand{\thesubsection}{\arabic{subsection}}
%\newtheorem{thm}{Theorem}[section]
%\renewcommand{\thethm}{\arabic{subsection}.\arabic{thm}}

\usepackage{thmtools}
\usepackage{hyperref}
\hypersetup{
    colorlinks=true,
    linkcolor=blue,
    filecolor=magenta,      
    urlcolor=cyan,
    pdftitle={Overleaf Example},
    %pdfpagemode=FullScreen,
    }
%\usepackage{stix-type1}%\turnediota
\theoremstyle{definition}

\newtheorem{constructor}{Constructor}[section]
\newtheorem{iaxiom}{Axiom}
\newtheorem{axiom}{Axiom}
\newtheorem*{axiom*}{Axiom}
\newtheorem{definition}{Definition}[section]
\newtheorem*{definition*}{Definition}
\newtheorem{remark}{Remark}[section]
\newtheorem*{remark*}{Remark}
\newtheorem{recall}{Recall}[section]
\newtheorem*{recall*}{Recall}

\newtheorem{lemma}{Lemma}[section]
\newtheorem{theorem}{Theorem}[section]
\newtheorem*{theorem*}{Theorem}
\newtheorem{note}{Note}[section]
\usepackage{etoolbox}

%\newcommand*{\@Dashv}[2]{%
%  \reflectbox{$\m@th#1#2$}%
%}
\newcommand*{\separation}[3]{% 1)binder, 2)set, 3)condition
  \{#1\in #2: #3\}
}

% https://www.freecodecamp.org/news/simple-regex-tricks-for-beginners-3acb3fa257cb/
%  \\mathcal\{T\} |---> \topology
%  \\mathcal\{([^\}]*)\}   |--->  $1
\newcommand{\AND}{\ \ \&\ \ }
\newcommand{\any}{\mathrm{\ any}}
\newcommand{\borel}{\mathcal{B}}
\newcommand{\can}{\mathrm{can}}
\newcommand{\canonize}{\mathrm{canonize}}
\newcommand{\closure}{\mbox{cl}}
\newcommand{\complex}{\mathbb{C}}
\newcommand{\converges}{\longrightarrow}
\newcommand{\countable}{\mbox{\ countable}}
\newcommand{\ldefined}{\mbox{\ defined}}
\newcommand{\dom}{\mbox{dom}}
\newcommand{\Dom}{\mbox{Dom}}
\newcommand{\Done}{\noindent \textbf{Done.}\\}
\newcommand{\lexists}{\mbox{\ exists}}
\newcommand{\lforall}{\mbox{\ forall}}
\newcommand{\lunique}{\mbox{\ unique}}
\newcommand{\family}{\mathrm{\ family}}
\newcommand{\Fm}{\mathrm{Fm}}
\newcommand{\FuncSymb}{\mathrm{FuncSymb}}
\newcommand{\function}{\mathrm{\ function}}
\newcommand{\functional}{\mbox{\ functional}}
\newcommand{\IFF}{\ \Longleftrightarrow\ }
\renewcommand{\iff}{\ \longleftrightarrow\ }
\newcommand{\IMPLIES}{\ \Longrightarrow\ }
\renewcommand{\implies}{\ \longrightarrow\ }
\newcommand{\Inf}{\mathrm{Inf}}
\newcommand{\injective}{\mbox{\ injective}}
\newcommand{\injection}{\mbox{\ injection}}
\newcommand{\interior}{\mbox{int}}
\newcommand{\intersection}{\bigcap}
\newcommand{\isempty}{\mathrm{\ empty}}
\newcommand{\LB}{\mathrm{LB}}
\newcommand{\nghbset}{\mathcal{N}}
\newcommand{\nonempty}{\mathrm{\ nonempty}}
\newcommand{\Ord}{\mathrm{Ord}}
\newcommand{\power}{\mathcal{P}}
\newcommand{\pr}{\mathrm{pr}}
\newcommand{\fst}{\mathrm{fst}}
\newcommand{\snd}{\mathrm{snd}}
\newcommand{\PredSymb}{\mathrm{PredSymb}}
\newcommand{\Proof}{\noindent \textbf{Proof.}\\}
\newcommand{\ran}{\mathrm{ran}}
\newcommand{\Ran}{\mathrm{Ran}}
\newcommand{\real}{\mathbb{R}}
\newcommand{\relation}{\mathrm{\ relation}}
\newcommand{\set}{\mathrm{\ set}}
\newcommand{\singleton}{\mathrm{\ singleton}}
\newcommand{\semantics}{\mathrm{semantics}}
\newcommand{\semialgebra}{\mathcal{S}}
\newcommand{\Sup}{\mathrm{Sup}}
\newcommand{\surjective}{\mbox{\ surjective\ on\ }}
\newcommand{\Term}{\mathrm{Term}}
\newcommand{\Type}{\mathrm{Type}}
\newcommand{\topology}{\mathcal{T}}
\newcommand{\total}{\mbox{\ total\ on\ }}
\newcommand{\transitive}{\mathrm{\ transitive}}
\newcommand{\union}{\bigcup}
\newcommand{\lTrue}{\mathrm{True}}
\newcommand{\lFalse}{\mathrm{False}}
\newcommand{\lUncertain}{\mathrm{Uncertain}}
\newcommand{\lif}{\mathrm{if}}
\newcommand{\lthen}{\mathrm{then}}
\newcommand{\lelse}{\mathrm{else}}
\newcommand{\lexception}{\mathrm{exception}}
\newcommand{\lundef}{\mathrm{undef}}
\newcommand{\univocal}{\mathrm{\ univocal}}
\newcommand{\univocalat}{\mathrm{\ univocal\ at\ }}
\newcommand{\UB}{\mathrm{UB}}
\newcommand{\Un}{\mathcal{U}}
\newcommand{\Var}{\mathrm{Var}}

\newcommand{\notvarphi}{\phi}
\renewcommand{\phi}{\varphi}

\newcommand{\comment}[1]{\quad\mbox{(#1)}}
% local commands
%\preto\section{\undef\blipo} %etoolbox
\preto\subsection{\undef\blipo} %etoolbox
% https://tex.stackexchange.com/questions/105684/limit-macro-newcommand-scope-to-section
\usepackage{graphicx}
\usepackage{authblk}

\setlength\parindent{0pt} % whole doc of \noindent
\usepackage{soul}
%\usepackage{mathgl}

\fi % guard
\ifdefined\topic
\else
\def\topic{No Titel}
\fi
\title{\topic}
\author[1]{Georgy Dunaev}
\begin{document}
\maketitle
\fi

\section{Value} % of a function.
Definition of the value of the function: (in the context of 3-valued logic)

\begin{definition}(quasi-equality)
\[
(s\simeq t) :\IFF (((s \ldefined)\lor(t \ldefined)) \implies (s=t))
\]
\end{definition}

\begin{definition}(undefined class)
\[
\lundef :\simeq (\varepsilon x: \bot)
\]
\end{definition}

\begin{definition}(if-then-else)todo
\[
(\lif\ \varphi\ \lthen\ t_1\ \lelse\ t_2) :\simeq (\{z|((z\in t_1)\land\varphi)\lor((z\in t_2)\land\neg\varphi)\})
\]
Side conditions:\\
($z$ is not free in $t_1$),\\
($z$ is not free in $t_2$),\\
($z$ is not free in $\varphi$).\\
\end{definition}
What if $\varphi$ is assigned with $U$ (i.e. it is assigned with the "uncertain" truth value)?\\
If $(I(\varphi)=U)$, then 
\[
(I((\lif\ \varphi\ \lthen\ t_1\ \lelse\ t_2))=I(\lundef))
\]
should follow.\\
Btw, $I(\lundef)=\emptyset$.\\

Note 1:\\
Some things are undefined because of "plural" interpretation, some because of "nullar" interpretation. Only those which have "singular" interpretation are defined.\\
Note 2:\\
Technically almost everywhere where we use "defined" we mean well-defined, but write just "defined" for shortness.\\
Note 3:\\
If we explicitly assume the global choice, then epsilon operator may be assumed to be defined.\\
If we didn't, epsilon operator is can still be defined, but not well-defined.\\

Now, let's talk about how to understand Hilbert's definition of quantifiers via predicate and epsilon operater: (note the two different epsilon operators!)
\[
((\varepsilon x:\varphi) \lunique) \IFF (! x:\varphi)
\quad\quad\quad\quad\quad\quad\quad\quad\quad\quad
\]
%:\IFF \varphi[(\epsilon x:\varphi)/x]
%\]
\[
((\varepsilon x:\varphi) \lexists) \IFF (\exists x:\varphi):\IFF \varphi[(\epsilon x:\varphi)/x]
\]
\[
(\neg((\varepsilon x:\neg\varphi) \lexists)) \IFF (\forall x:\varphi):\IFF \varphi[(\epsilon x:\neg\varphi)/x]
\]
If $(\varepsilon x:\varphi)$ is not well-defined for particular $\varpi$, then the whole predicate $\varphi[(\varepsilon x:\varphi)/x]$ is interpreted, in general, as an "uncertain".\\
Therefore for such definition we need a well-defined(total) variant. (extended with some default value)
\\
Btw,
\[
((\varepsilon x:\varphi) \lforall) :\IFF (\neg((\varepsilon x:\neg\varphi) \lexists)).
\]
For all connectives, if there is a U value, then the otcome will also be U.\\
Including equivalences!\\
Therefore "definition" is working correctly, only in the case, when $(\varepsilon x:\varphi)$ is well-defined.\\
We can distinguisch $\epsilon$ which is total, and $\varepsilon$ which is not defined on the empty predicates.\\
Frankly speaking, since it assumes the uniqueness of the value, we need to talk about\\

There are 3 natural ways of dealing with a default value of the choice.\\
1) $\tau^\emptyset$. Here no value is given. It is unique, but it doesn't exist.\\
2) $\tau^{\exists!}$. Here one particular value is set as a default. It may even be specified: e.g. $\tau^{\exists!}_\emptyset$, or $\tau^{\exists!}_\Un$ or, in general, $\tau^{\exists!}_K$, which is the same as even better: $\tau^{\{K\}}$.\\
3) $\tau^\Un$. Here any class from the conglomerate of all classes $\Un$ value is a potential default value. It exists, but it isn't unique.\\


Our convention:\\
$\tau := \tau^\Un$\\
$\epsilon := \tau^{\exists!}_\emptyset = \tau^{\{\emptyset\}}$\\
$\varepsilon := \tau^\emptyset$\\
In usual mathematics, usually either $\tau^{\exists!}$ or $\epsilon$ are used.\\

%1) $t$ -- total choice functional symbol ($\tau$, since it reminds both bourbaki notation and starting letter of the word "total");\\
%2) $\tau$ -- partial choice functional symbol ($\tau$, since it reminds both bourbaki notation and starting letter of the word "total");\\

%2) $c$ -- choice functional symbol (it works also on proper classes. It is not a strict extension of usual global choice, since from global choice proper classes are bijective with classs of all sets, and it has an element, let's say the empty set. But, is it necessary to choose th choice fuction? I don't know.);\\
%3) total $\epsilon$-operator, which returns empty set in case of definable everywhere false predicates is given as an argument;\\
%4) partial $\varepsilon$-operator, which is undefined on everywhere false definable predicates.\\

Note 1: operators work only on definable conglomerates.\\
Note 2: I think predicate is a little bit more general than the bounded formula since bounded formulas are corresponding to definable conglomerates.\\
 
For Hilbert's variant of the definition of quantifiers, it is handy to use \#3: the total $\epsilon$-operator.
 
%$\bdbvbH$
%\symbol{216}

%$\symbol{208}$
%qqqq
\newcommand{\uncert}{\mbox{\symbol{208}}}
%$\uncert$
%ggggg
To denote the constructor of an uncertain formula, we will use the symbol $\uncert$, because it is symmetrical (neither $\bot$, nor $\top$ are  symmetrical ). And it looks like crossed D, and "defined" starts with D, and it is connected with not being defined: normally(but not always!), if $f$ is a functional symbol, $f(...,\lundef,...)$.\\
\[
I_\Term(\uncert) \coloneqq U
\]
\pagebreak


We have not always well-defined epsilon operator.\\
(tau operator is always well-defined, epsilon operator is well-defined only on bounded formulas which specify "singletone-conglomerates"(if we could collect classes).\\

Here we come to the point, where we nee more relaxed notion of truthness, to include $P((\varepsilon X: \varphi))$, when predicate P is true for all classes "X" satisfying $\varphi$.

It is interesting, that quasi-equality is naturally arises as... oh no!
We need to introduce a mix "exists and forall" quantifier and semantics, so the quasi-equality will natuarally arise.
I mean the following:
\begin{itemize}
\item
$P(t_1,...,t_n)=\lTrue$, iff two conditions are satisfied:\\
1) $I_\Term(t_j)$ is non-empty\\
2) for all sequences of object $(q)_{j=1}^n$ , such that for all j from 1 to n, $q_j\in I_\Term(t_j)$: $I_{\PredSymb}(P)(q_1,...,q_n)=\lTrue$.
\item $P(t_1,...,t_n)=\lFalse$, iff two conditions are satisfied:\\
1) $I_\Term(t_j)$ is non-empty\\
2) for all sequences of object $(q)_{j=1}^n$ , such that for all j from 1 to n, $q_j\in I_\Term(t_j)$: $I_{\PredSymb}(P)(q_1,...,q_n)=\lFalse$.
\item
$P(t_1,...,t_n)=\lUncertain$, in all other cases.\\
For example, if there exists $j\in\{1,...,n\}$, such that $I_\Term(t_j)$ is empty;\\
or if neither of the second requirements hold,\\
i.e. if the function $\left((q)_{j=1}^n \mapsto I_{\PredSymb}(P)(q_1,...,q_n)\right)$ is not a constant on $I_\Term(t_1)\times\dots\times I_\Term(t_n)$.
\end{itemize}


Assume $a$, $b$ are well-defined sets.\\
Then $
(\tau x: x \in \{a,b\}) 
$ is well-defined, since using it assumes existence of a choice conglomerate-function.\\
Meanwhile,
$
(\varepsilon x: x \in \{a,b\}) 
$
 is undefined. But it is not (semantically) equal to $\lundef$.\\
It is neither equal, nor not equal in theory, because, some(in fact both) sides are not well-defined.\\
But they are quasi-equal (syntactically), since they are both not well-defined.\\
 
Remark: Here by semantical equality i mean the equaly of all interpretations of two term under any assignment of values.

Remark: actually we can use 6 truth values.\\
 True = (present, stable(True))\\
False = (present, stable(False))\\
U1    = ( absent, stable(True))\\
U2    = ( absent, stable(False))\\
U3    = (present, unstable)\\
U4    = ( absent, unstable)\\
Connectives are contagious in the first argument, it means the following:\\
Let's say \# is an interpretation of a connective, and v is a truth value, then\\
 (absent, ...) \# v = (absent, ...)\\
 (note, that "..." in lhs and rhs are not necessary the same)\\
 (..., stable(True)) $\land$ (..., stable(True)) =  (..., stable(True))\\
But two predicates can stabilize each other, so these truth values does not have a truth table, that's why we will not use them.\\

Also therefore Uncertain OR Uncertain is not necessary Uncertain in such semantics!
Therefore we need to require even more: the certainity of every subformula, otherwise it will complicate the system too much?\\

e.g. $(f(x)=5) \lor (x\notin\dom(f))$.\\
Or will it require too much to require?\\
Yes, it will be too much, since then\\
$(f\function)\land(x\in\dom(f)) \implies (f(x)=f(x))$\\
will not work normally.\\
Therefore we need something more interesting like:\\
 $(\lFalse\implies\lUncertain) := \lTrue$.\\
Or even brave \& better: reject the whole idea of having a defined truth table for uncertain values(or at least having a full/complete truth table).
Such idea was wrong, because it is not true even in completely classical logic of existence: the formula "$(x\in y)$" is not valid, so it has no any particular defined truth value.\\

Assigment maps each variable to some singletone.\\
Let's come back to:
$(f\family)\land(x\in\dom(f)) \implies (f(x)=f(x))$\\
For the case when something goes wrong, "(f(x)=f(x))" is not true, since it being true assumes both-hand-sides to be defined.


Recall:\\
1) The term $(\varepsilon x: x \in \{a,b\})$ is not unique.\\
2) The formula $a=(\varepsilon x: x \in \{a,b\})$ is not stable under differen assignment.\\
3) Something of the following gives impression of something that should be true:
$(a=(\varepsilon x: x \in \{a,b\}))\lor(b=(\varepsilon x: x \in \{a,b\}))$(bad)\\
$z=(\varepsilon x: x \in \{a,b\}) \implies (a=z)\lor(b=z)$(bad)\\
$z\simeq (\varepsilon x: x \in \{a,b\}) \implies (a=z)\lor(b=z)$(bad)\\
$t\simeq (\varepsilon x: x \in \{a,b\}) \implies (a=t)\lor(b=t)$(bad)\\
$t:\simeq (\varepsilon x: x \in \{a,b\}) \implies (a=t)\lor(b=t)$(bad)\\
$t:\simeq (\varepsilon x: x \in \{a,b\}) \implies (a\simeq t)\lor(b\simeq t)$(bad)\\
We need to determine what is valid and what is not.\\
Resolution:\\
We should not use variables for definitions, because definitions are not always well-defined. And variables are always well-defined.\\
We must use new constants instead.\\
It is a more delicate mechanism than just the use of quasi-equality(as in W.Farmer's article).\\
So we have a constant $c:=(\varepsilon x: x \in \{a,b\})$.
But still $(a\simeq c)\lor(b\simeq c)$ is false.
Maybe $(\exists q\in\{a,b\}:c=q)$? No.(YES or NO? i'm not sure)

In order for lemmas from (link shen-logic-part2-2.pdf , page 137) to work, we need assign non-sigleton sets to variables.\\
But we want to avoid tedious assumtions like the following:
$(x_1\ldefined)\land...\land(x_n\ldefined)\implies \varphi$.\\
Therefore for theorems on the upper level assigment are considered to be with singletones only, but within the evaluation variables may be assigned to a set with non-unique values.\\
The previous comment about introducing constants will patch the need for shortening of the statement while it is not possible to do via variable.\\
But what about the cases when it should be a term, depending on some parameters?\\
Resolution: we use only terms, and this mechanism of definitions of terms is outside the logic of the system.\\
Or we can introduce $:\equiv$ somehow? YES.\\%$:\approx$
\begin{definition}(interpretation of definitional equality)\\
\st{$I^v_\Fm(A\equiv B):=(I^v_\Term(A) = I^v_\Term(B))$}.NO! WRONG!\\
$I^v_\Fm(A\equiv B):=\forall(w\ \mbox{evaluation}):(I^w_\Term(A) = I^w_\Term(B))$.\\
Will such definition work with a variable? I highly doubt. 
So it's better to keep such definitions of undefined values outside of the logic until there will be an explicit necessity for them. (Probably there will be no.)
But this definition works well for defining functional symbols, and, in particular, constants.
\end{definition}

Some statements:\\
1) VALID, TRUE: $(\varepsilon x: x \in \{a,b\})\in \{a,b\}$.\\
2) INVALID: $(\varepsilon x: x \in \emptyset) \in \emptyset$. It is neither true, nor false, because the interpretation of the first term is an empty conglomerate of classes.\\
3) INVALID: $(\varepsilon x: x \in \emptyset) \notin \emptyset$. It is neither true, nor false, because the interpretation of the first term is an empty conglomerate of classes.\\
4) VALID, TRUE: $(x\in y) \implies (x\in y)$.

============================================\\
Idea: Truth values of predicates are subconglomerates of $(2\times)\Un\times\dots\times\Un$ of values, where $I(P)$ is true.\\
"(Always)True" corresponds to the whole subconglomerate $(1\times)\Un\times\dots\times\Un$.\\
"(Always)False" corresponds to the empty subconglomerate $(0\times)\emptyset\times\dots\times\emptyset$, which is simply $\emptyset$.\\
"Uncertain" corresponds to any other subconglomerate.\\
All products here are given $n$ times: $n$ is a valency of the predicate symbol $P$.
============================================\\

For example, the following predicates:\\
\begin{itemize}
\item (of valency 0):\\
"$\bot$" corresponds to 0.("False")\\
"$\top$" corresponds to 1.("True")\\
"$\uncert$" corresponds to 2.("Uncertain")\\
(from now on, we will forget about "$2\times$" form the definition, for simplicity)
\item (of valency 1):\\
"(X unique)" corresponds to the collection of all conglomerate-semisingletones.\\
"(X exists)" corresponds to the collection of all nonempty conglomerates.\\
"(X (well-)defined)" corresponds to the collection of all nonempty conglomerate-semisingletones, i.e. to the collection of all conglomerate-singletones\\
"empty" \ corresponds to the collection of empty conglomerates, i.e. just $\{\emptyset\}$.
\item (of valency 2):\\
"$(x\in y)$" corresponds to collection of all ordered pairs of nonempty conglomerates, where each class from first conglomerate is in each class of the second conglomerate.\\
"$(x\notin y)$" corresponds to collection of all ordered pairs of nonempty conglomerates, where each class from first conglomerate is not in any class of the second conglomerate.
\end{itemize}
\begin{definition}(Theorem)
Theorem is a predicate which has the truth value "True".\\
\end{definition}
%$\tau t$
%Therefore.
%This is something unresolved, we need some time...


% and define the Uncertain value, like:\\
%Uncertain = Non-existent OR Non-constant\\
%\,\quad\quad = Absent OR Unstable\\
%\,\quad\quad = Void OR Varying\\
%But then, how connectives will work? They will be "contagious" take "the worst case"? %Not necessary: two predicates can stabilize each other.\\
Idea 2:\\

Actually, variables are ordered: $x_1,x_2,x_3,\dots$.
So we can consider the following correspondence evaluation of the formula $x_2\in x_4$ via assignment function $w$, with 
\[
(1\times)\Un\times \{w(x_2)\}\times\Un\times\{w(x_4)\}\times\Un\times\Un\times\Un\dots,
\]
where "infinite product" is abuse of notation, which means a collection of functions $f$ from $\mathbb{N}_0$ to $\power\Un$, such that $f(0)=1$. %\subseteq 

\iffalse
\symbol{200}
\symbol{201}
\symbol{202}
\symbol{203}
\symbol{204}
\symbol{205}
\symbol{206}
\symbol{207}
\symbol{208}
\symbol{209}
\symbol{210}
\symbol{211}
\symbol{212}
\symbol{213}
\symbol{214}
\symbol{215}
\symbol{216}
\symbol{217}
\symbol{218}
\symbol{219}\\
\symbol{220}
\symbol{221}
\symbol{222}
\symbol{223}
\symbol{224}
\symbol{225}
\symbol{226}
\symbol{227}
\symbol{228}
\symbol{229}
\symbol{230}
\symbol{231}
\symbol{232}
\symbol{233}
\symbol{234}
\symbol{235}
\symbol{236}
\symbol{237}
\symbol{238}
\symbol{239}\\
\symbol{240}
\symbol{241}
\symbol{242}
\symbol{243}
\symbol{244}
\symbol{245}
\symbol{246}
\symbol{247}
\symbol{248}
\symbol{249}
\symbol{250}
\symbol{251}
\symbol{252}
\symbol{253}
\symbol{254}
\symbol{255}
\fi

%\symbol{253A}
%not everything that is undefined is quasi.


\begin{definition}(value)
\[
F(x) :\simeq \left(\lif (F\family)\land(x \in \dom(F))\ \lthen\ \union\{s|(x,s)\in F\}\ \lelse\ \lundef\ \lexception\_\lelse\ \lundef\right)
\]
\end{definition}


\section{Some simple theorems about classes}
\begin{definition}(some syntactic sugar)
\[
(x\ \mbox{any}) :\iff \top
\]
Context: (x:Var)
%Side condition:\\
%Term $x$ is a variable.
\end{definition}

\begin{definition}(Semi-singleton)
\[
(s\ \mbox{ss}) :\iff (!y:y\in s)
\]
Side condition:\\
Variable $y$ is not in free variables of term $s$.
\end{definition}

\begin{definition}(Semi-singleton subset relation)
\[
(s\ \mbox{sss}\ t) :\iff ((s\ \mbox{ss})\land(s\subseteq t))
\]
\end{definition}

\begin{definition}(Function constructor)
\[
(\lambda x: t) := \{(x,t)|(x\ \mbox{any})\}
\]
Context: ($x$:Var) ($t$:Term)\\
Comment: Term $t$ can contain $x$ as a free variable.\\
Side conditions:\\ 
(variable $x$ is bound by class collection in r.h.s)\\
(variable $x$ is bound by lambda abstraction in l.h.s)\\
%No side conditions.
\end{definition}

\begin{definition}(Family constructor)
\[
(\Lambda x: t) := \{(x,s)|(s\ \mbox{sss}\ t), x\ \mbox{any}\}
\]
% s is a bound variable
Term $t$ can contain $x$ as a free variable.\\
Side conditions:\\
Variables x and s are different.\\ %it should be a normal convention for any definition.
Variable $s$ is not in free variables of $t$.\\
As term $s$ can contain $x$ as a free variable.\\
\end{definition}

\begin{definition*}(LOCAL definition of a value of a function, value of a family, for the next theorem)
\[
f(x) := \bigcup \{s|(x,s)\in f\}
\]
\end{definition*}
% \coloneqq
% \vcentcolon = 

\begin{theorem}(Evaluation of a function)
Context:\\
(x:Var)(t:Term)(z:Term)
\[
(z\set) \land (t[z/x]\set) \implies
\]
\[
(\lambda x: t)(z) = t[z/x]
\]
\end{theorem}
\Proof
To prove the equality, we use set extensionality and MP.
\[
((\lambda x: t)(z) \subseteq t[z/x] )
\land
((\lambda x: t)(z) \supseteq t[z/x] )
\implies
((\lambda x: t)(z) = t[z/x] )
\]
($\subseteq$)\\
Assume 
\[
q\in (\lambda x: t)(z)
\]
By definition of $-(-)$ (evaluation),
\[
q\in\union \{b|(z,b)\in(\lambda x:t)\}.
\]
By destruction of the union,
\[
\exists s\in \{b|(z,b)\in(\lambda x:t)\}: q\in s.
\]
Note for the following step: since variable b has no free occurences in t, $t[s/b]=t$.
\[
s\in \{b|(z,b)\in(\lambda x:t)\}
\IMPLIES
(s\set) \& ((z,s)\in(\lambda x:t))
\]
By definition of $\lambda$,
\[
(z,s)\in(\lambda x:t)
\IMPLIES
(z,s)\in \{(x,t)|\top\}
\]
Unfolding the notation:
\[
(z,s)\in \{(x,t)|\top\}
\IMPLIES
(z,s)\in \{w|\exists x: w = (x,t) \land \top\}
\]

\[
(z,s)\in \{w|\exists x: w = (x,t) \land \top\}
\IMPLIES
\exists x: (z,s) = (x,t) \land \top
\]
Notes: w is free in t, w is different from x. Since otherwise we would write $(x[(z,s)/w],t[(z,s)/w])$.

Take such x. By defining property of ordered pairs,
\[
(z,s) = (x,t)
\IMPLIES
(z=x) \& (s=t)
\]
\[
(q\in s)\& (s=t)
\IMPLIES
q\in t
\]
\[
z=x
\IMPLIES
t[z/x]=t
\]
Therefore
\[
(q\in t) \& (t[z/x]=t)
\IMPLIES
(q \in t[z/x])
\]
\ \\
\ \\
($\supseteq$)\\
Assume $(q \in t[z/x])$.\\
Let $s:=t[z/x]$. Then $(q\in s)$.\\
Goal1:$(z,s)\in(\lambda x:t)$.\\
\[
x[z/x]=z
\]
\[
(z,t[z/x]) = (x[z/x],t[z/x])=(x,t)[z/x]
\]
Goal1: $(x,t)[z/x]\in\{(x,t)|(x\mbox{ any})\}$\\
Goal2: $\exists x:(x,t)=((x,t)[z/x])$\\
x:=z\\
There are no free occurencies of $x$ in $((x,t)[z/x])$, because of the substitution.\\
$c:=((x,t)[z/x])$\\
\[
(((x,t)[z/x])=c)\implies\exists x: (x,t)=c
\]
Goal' is done.\\

From the premise, $(x,t)[z/x]$ is a set.\\
Therefore Goal1 is proven.\\
So $(x,t)[z/x]\in (\lambda x:t)$.\\
So $(z,t[z/x])\in (\lambda x:t)$.\\
So $t[z/x]\in \{b|(z,b)\in(\lambda x:t)\}$.\\
Since $(q \in t[z/x])$, we also have
\[
q\in \bigcup \{b|(z,b)\in(\lambda x: t)\}
\]
So $q\in (\lambda x: t)(z)$.\\
And we are done.\\
\Done


\begin{theorem}(Evaluation of family)
\[
(\Lambda x: t)(z) = t[z/x],
\]
and in particular:
\[
(\Lambda x: t)(x) = t.
\]
%Context:
%(q term)
\end{theorem}
\Proof
By definitions,
\[
(\Lambda x: t)(z) =.
\]


\Done

\subsection{Lambda operatorOLD}
\begin{definition}
\label{def:lambda operator}not good anymore: doesn't work for classes
Lambda operator.
\[
(\lambda x: t) = \{p|\exists x: p = \langle x,t \rangle\}
\]
\end{definition}

\ifdefined\qwertyflag
\undef\qwertyflag
\end{document}
\fi